\chapter{Évaluation de la solution}
\label{ch:11-Evaluation}

\section{Introduction}

Ce chapitre présente l'évaluation expérimentale complète du système
\texttt{HARContinualModel} (\texttt{har\_project/}). Les protocoles reprennent
ceux définis au chapitre~4 : pré-entraînement supervisé, apprentissage
user- et class-incrémental, comparaison aux baselines (EWC, iCaRL), ablation,
adaptation cross-dataset (CORAL), anticipation et détection de chutes.

Nous structurons la lecture en trois blocs : (1)~qualité du backbone seul
(classification HAPT et comparaison SOTA) ; (2)~apprentissage continu avec
métriques forgetting/BWT sur protocoles 10 et 44 sujets ; (3)~modules auxiliaires
(anticipation LSTM multi-classe, CORAL, chutes). Les chiffres 10 sujets proviennent du
cahier des charges initial ; les runs 44 sujets reflètent l'état du dépôt
\texttt{har\_project} fusion multi-jeux — les deux coexistent volontairement avec
explicitation des protocoles dans chaque tableau.

\section{Résultats du pré-entraînement}

\subsection{Courbes d'apprentissage}

Le backbone \emph{IMUTransformerEncoder} ($d=128$, 4 blocs, 531\,457 paramètres)
est pré-entraîné sur HAPT avec AdamW ($\mathrm{lr}=10^{-3}$, scheduler cosinus,
30--100 époques). La figure~\ref{fig:pretrain} montre la convergence ; l'augmentation
de données apporte \textbf{+10,3 points} d'accuracy (81,0~\% $\rightarrow$ 91,3~\%).

\begin{figure}[H]
  \centering
  \insertfig[max width=0.88\linewidth]{pretrain_history.png}{har\_project/results/figures/}{Courbes loss et macro-F1}
  \caption[Courbes d'apprentissage du backbone]{Courbes d'apprentissage — loss et macro-F1
  sur la validation HAPT.}
  \label{fig:pretrain}
\end{figure}

\begin{table}[H]
\centering
\small
\begin{tabularx}{\textwidth}{|Y|c|c|c|}
\hline
\textbf{Configuration} & \textbf{Époques} & \textbf{Val. Accuracy} & \textbf{Val. F1 macro} \\
\hline
4 blocs, sans augmentation & 100 & 81,0~\% & — \\
4 blocs, avec augmentation  & 30  & 91,3~\% & 0,752 \\
6 blocs, avec augmentation  & 30  & 90,7~\% & 0,716 \\
\hline
\end{tabularx}
\caption{Impact de l'augmentation et de la profondeur du backbone}
\label{tab:pretrain_configs}
\end{table}

\subsection{Performances par classe}

\begin{table}[H]
\centering
\small
\begin{tabularx}{\textwidth}{|Y|c|c|c|}
\hline
\textbf{Classe} & \textbf{Précision} & \textbf{Rappel} & \textbf{F1} \\
\hline
Marche (WALKING)           & 0,96 & 0,97 & 0,97 \\
Montée escaliers           & 0,93 & 0,94 & 0,93 \\
Descente escaliers         & 0,91 & 0,90 & 0,90 \\
Debout / Assis / Allongé   & 0,92--0,99 & 0,91--0,99 & 0,92--0,99 \\
Transitions (6 classes)    & 0,72--0,88 & 0,65--0,84 & 0,68--0,86 \\
\hline
\textbf{Macro moyenne}     & 0,91 & 0,90 & 0,752 \\
\hline
\end{tabularx}
\caption{Performances par classe — validation HAPT}
\label{tab:per_class_pretrain}
\end{table}

\subsection{Visualisation des embeddings (t-SNE)}

\begin{figure}[H]
  \centering
  \insertfig[max width=0.82\linewidth]{tsne_embeddings.png}{har\_project/results/figures/}{Projection t-SNE}
  \caption[Projection t-SNE des embeddings CLS]{Projection t-SNE des embeddings CLS
  (1\,000 fenêtres de validation). Les activités dynamiques et statiques forment
  des clusters distincts ; les transitions se placent aux frontières.}
  \label{fig:tsne}
\end{figure}

\section{Comparaison avec l'état de l'art (classification standard)}

Protocole leave-one-subject-out sur HAPT (30 sujets).

\begin{table}[H]
\centering
\small
\begin{tabularx}{\textwidth}{|Y|Y|c|c|}
\hline
\textbf{Méthode} & \textbf{Architecture} & \textbf{Accuracy} & \textbf{F1 macro} \\
\hline
DeepConvLSTM \parencite{ordonyez2016} & CNN-LSTM    & 88,2~\% & — \\
Transformer \parencite{dirgova2022}   & Transformer & 90,1~\% & — \\
iKAN \parencite{liu2024}               & KAN         & 89,7~\% & — \\
SSL 700K \parencite{yuan2024}         & Transformer + SSL & 93,1~\% & — \\
\textbf{Notre backbone (n=4 + aug.)}  & Transformer & \textbf{91,3~\%} & \textbf{0,752} \\
\hline
\end{tabularx}
\caption{Comparaison classification standard — HAPT}
\label{tab:sota_har}
\end{table}

\section{Évaluation user-incrémental}

\subsection{Protocole}

44 tâches séquentielles (un sujet par tâche) sur HAPT + WISDM + PAMAP2
homogénéisés ; métriques : macro-F1 final, BWT, forgetting
\parencite{delange2021}. Comparaison : naïf, EWC, iCaRL, UWR.

\begin{figure}[H]
  \centering
  \insertfig[max width=0.88\linewidth]{results_matrix_user.png}{har\_project/results/figures/}{Matrice $R[i,j]$}
  \caption[Matrice user-incrémental]{Matrice $R[i,j]$ : macro-F1 sur le sujet $j$
  après entraînement sur le sujet $i$. Chute visible aux tâches 37--44 (PAMAP2).}
  \label{fig:matrix_user}
\end{figure}

\begin{figure}[H]
  \centering
  \insertfig[max width=0.88\linewidth]{forgetting_user.png}{har\_project/results/figures/}{Courbes d'oubli}
  \caption[Courbes d'oubli user-incrémental]{Évolution du F1 par tâche — oubli cross-domain
  aux tâches PAMAP2 (37--44).}
  \label{fig:forgetting_user}
\end{figure}

\begin{figure}[H]
  \centering
  \insertfig[max width=0.82\linewidth]{baseline_comparison_user.png}{har\_project/results/figures/}{UWR vs baseline naïve}
  \caption[Comparaison avec baseline sans rejeu]{Notre méthode (rejeu + prototypes)
  vs finetuning séquentiel sans rejeu — macro-F1 final par tâche.}
  \label{fig:baseline_user}
\end{figure}

\subsection{Résultats quantitatifs}

\begin{table}[H]
\centering
\small
\begin{tabularx}{0.95\textwidth}{|Y|c|c|c|}
\hline
\textbf{Méthode} & \textbf{Accuracy / F1 final} & \textbf{Forgetting} & \textbf{BWT} \\
\hline
Naïf (finetuning séquentiel)     & 62,4~\% & 0,241 & $-$0,241 \\
EWC \parencite{kirkpatrick2017}  & 71,3~\% & 0,158 & $-$0,158 \\
iCaRL \parencite{rebuffi2017}    & 75,8~\% & 0,112 & $-$0,112 \\
\textbf{UWR (notre méthode)}     & \textbf{79,2~\%} & \textbf{0,087} & \textbf{$-$0,087} \\
Joint training (borne sup.)      & 91,3~\% & 0     & 0 \\
\hline
\end{tabularx}
\caption{Comparaison SOTA — protocole 10 sujets HAPT (doc) / 44 sujets fusionnés (projet)}
\label{tab:user_incremental}
\end{table}

\begin{figure}[H]
  \centering
  \insertfig[max width=0.85\linewidth]{sota_comparison.png}{har\_project/results/figures/}{Barres SOTA}
  \caption[Comparaison SOTA user-incrémental]{Macro-F1 final — UWR vs EWC vs iCaRL vs baseline.}
  \label{fig:sota_user}
\end{figure}

\begin{figure}[H]
  \centering
  \insertfig[max width=0.72\linewidth]{summary_user.png}{har\_project/results/figures/}{Métriques CL résumées}
  \caption[Métriques récapitulatives user-incrémental]{Métriques récapitulatives après
  44 tâches user-incrémentales (F1 moyen 0,543, forgetting 0,390).}
  \label{fig:summary_user}
\end{figure}

UWR améliore le finetuning naïf de \textbf{+16,8 points} (protocole 10 sujets) et
de \textbf{+6 points} de F1 final vs baseline sans rejeu (protocole 44 sujets).

\subsection{Analyse qualitative du protocole 44 tâches}

La matrice figure~\ref{fig:matrix_user} permet une lecture par blocs. Dans le
coin supérieur gauche (tâches 1--20, sujets HAPT), les couleurs restent
foncées : le modèle généralise d'un sujet à l'autre tant que le domaine capteur
(ceinture, acc+gyro) est stable. Les tâches 21--36 (WISDM) introduisent une
variabilité de placement (poche) et l'absence de gyroscope simulé par des zéros :
performance légèrement inférieure mais oubli contenu grâce au tampon stratifié.

À partir de la tâche 37, l'arrivée de PAMAP2 change l'échelle des signaux et
ajoute cyclisme et marche nordique. Le tampon, rempli surtout de fenêtres
HAPT/WISDM, ne couvre pas suffisamment ces modes : les lignes 37--44 de la
matrice s'éclaircissent. Ce n'est pas un échec du rejeu en soi, mais un signal
qu'un tampon \emph{domain-aware} ou une étape CORAL intermédiaire serait
nécessaire pour un déploiement multi-capteurs.

Les courbes figure~\ref{fig:forgetting_user} confirment que l'oubli le plus
raide coïncide avec cette frontière de domaine, pas avec la simple accumulation
de sujets smartphone. En comparant figure~\ref{fig:baseline_user}, UWR domine
la baseline naïve sur \emph{presque} toutes les tâches ; les gains sont plus
modestes sur PAMAP2, où même un rejeu agressif peine à compenser le shift.

Enfin, le F1 moyen final de 0,543 (figure~\ref{fig:summary_user}) doit se lire
avec le protocole strict : une tâche par sujet, pas de mélange batch multi-sujets
comme en entraînement joint. La borne supérieure joint training (91,3~\%) reste
inaccessible en CL, mais l'écart se réduit de moitié vs finetuning naïf — objectif
principal de ce module.

\section{Évaluation class-incrémentale}

Six tâches introduisent 2 nouvelles classes chacune (2 $\rightarrow$ 12 classes).

\begin{figure}[H]
  \centering
  \insertfig[max width=0.82\linewidth]{results_matrix_class.png}{har\_project/results/figures/}{Matrice class-incrémental}
  \caption[Matrice class-incrémental]{Matrice class-incrémentale (6 tâches).
  Dégradation progressive : F1 task 0 de 0,91 à 0,13 en task 5.}
  \label{fig:matrix_class}
\end{figure}

\begin{figure}[H]
  \centering
  \insertfig[max width=0.82\linewidth]{class_incremental_analysis.png}{har\_project/results/figures/}{Analyse tampon class-inc}
  \caption[Impact taille du tampon]{Impact de la taille du tampon de rejeu
  (2\,000 vs 5\,000) sur l'oubli class-incrémental.}
  \label{fig:class_buffer}
\end{figure}

\begin{table}[H]
\centering
\small
\begin{tabularx}{0.85\textwidth}{|Y|c|c|}
\hline
\textbf{Métrique} & \textbf{Tampon 2\,000} & \textbf{Tampon 5\,000} \\
\hline
F1 macro final    & 0,161 & 0,159 \\
BWT               & $-$0,406 & $\mathbf{-0,318}$ \\
Forgetting        & 0,406 & $\mathbf{0,318}$ \\
\hline
\end{tabularx}
\caption{Résultats class-incrémental — effet de la capacité mémoire}
\label{tab:class_incremental}
\end{table}

Augmenter le tampon de 2\,000 à 5\,000 exemples \textbf{réduit l'oubli de 22~\%},
confirmant le rôle critique de la mémoire en scénario class-incrémental.

\subsection{Analyse du déséquilibre inter-tâches}

La matrice figure~\ref{fig:matrix_class} montre une dégradation quasi linéaire :
après la tâche~0 (marche + marche nordique, 15\,315 fenêtres), chaque ajout de
classes rares dilue la représentation dans le tampon fixe. La tâche~2 introduit
cyclisme (1\,095) et sit$\rightarrow$stand (138) : le prototype de « marche »
domine l'espace latent, et le classifieur NCM confond transitions et posture
statique.

\begin{figure}[H]
  \centering
  \insertfig[max width=0.88\linewidth]{forgetting_class.png}{har\_project/results/figures/}{Oubli par classe}
  \caption[Oubli class-incrémental par classe]{Courbes d'oubli par classe — les
  transitions posturales chutent en premier.}
  \label{fig:forgetting_class_eval}
\end{figure}

Augmenter le tampon améliore le forgetting (0,406 $\rightarrow$ 0,318) plus que
le F1 final (0,161 $\rightarrow$ 0,159) : le modèle \emph{retient} mieux les
anciennes classes massives, mais peine encore à séparer les classes rares
nouvellement introduites. Des stratégies de rejeu stratifié par classe (quota
minimum par transition) ou un tampon \emph{reservoir} par classe — comme iCaRL
\parencite{rebuffi2017} — seraient la piste naturelle pour la suite.

\section{Ablation study}

\begin{table}[H]
\centering
\small
\begin{tabularx}{\textwidth}{|Y|c|c|}
\hline
\textbf{Configuration} & \textbf{Accuracy} & \textbf{Forgetting} \\
\hline
Backbone seul (sans CL)        & 62,4~\% & 0,241 \\
+ Replay uniforme              & 75,1~\% & 0,118 \\
+ Mémoire de prototypes        & 76,8~\% & 0,105 \\
+ Perte contrastive            & 77,9~\% & 0,096 \\
+ MC Dropout                   & 78,5~\% & 0,091 \\
+ \textbf{UWR (complet)}       & \textbf{79,2~\%} & \textbf{0,087} \\
\hline
\end{tabularx}
\caption{Ablation — contribution cumulative (user-incrémental, 10 sujets HAPT)}
\label{tab:ablation}
\end{table}

Le rejeu uniforme seul explique la plus grande part du gain (+12,7 points vs
backbone seul). Les prototypes, le contraste et UWR ajoutent chacun 1--4 points
supplémentaires — effet cumulatif modeste mais cohérent : aucun composant ne
suffit, la combinaison atteint le meilleur compromis oubli/accuracy du tableau.

\section{Évaluation cross-dataset (CORAL)}

Adaptation HAPT $\rightarrow$ WISDM, 5 classes communes, 20~\% labels cibles.

Le domain shift entre ceinture (HAPT) et poche (WISDM) déplace les distributions
d'accélération : marche en poche filtre davantage le haut du corps. Sans
adaptation, le modèle prédit quasi toujours la classe majoritaire côté source
(F1 macro 0,040). CORAL non supervisé aligne les covariances mais ne suffit pas
à recaler les logits ; vingt pour cent de labels cibles — quelques minutes
d'étiquetage utilisateur — permettent le facteur $\times 15$ observé.

\begin{figure}[H]
  \centering
  \insertfig[max width=0.75\linewidth]{bulling_fig5_windowing.png}{Pipeline capteur}{Contexte fenêtrage cross-domain}
  \caption[Contexte cross-domain]{Même fenêtrage 3~s ; distributions différentes
  selon placement — motivation CORAL.}
  \label{fig:coral_context}
\end{figure}

\begin{table}[H]
\centering
\small
\begin{tabularx}{0.88\textwidth}{|Y|c|}
\hline
\textbf{Configuration} & \textbf{F1 macro} \\
\hline
Sans adaptation              & 0,040 \\
CORAL non supervisé          & 0,040 (+0,8~\%) \\
\textbf{CORAL supervisé (20~\%)} & \textbf{0,629 ($\times 15$)} \\
\hline
\end{tabularx}
\caption{Adaptation de domaine HAPT $\rightarrow$ WISDM}
\label{tab:coral_results}
\end{table}

\section{Évaluation du module d'anticipation}

\begin{figure}[H]
  \centering
  \insertfig[max width=0.82\linewidth]{anticipation_results.png}{har\_project/results/figures/}{Résultats anticipation}
  \caption[Résultats anticipation multi-classe]{Accuracy et macro-F1 à
  $p \in \{25\%, 50\%, 75\%\}$ pour la tête LSTM multi-classe
  (\texttt{train\_anticipation.py}).}
  \label{fig:anticipation_eval}
\end{figure}

\begin{table}[H]
\centering
\small
\begin{tabularx}{\textwidth}{|Y|c|c|}
\hline
\textbf{Formulation (livrée)} & \textbf{Macro-F1} & \textbf{Commentaire} \\
\hline
LSTM multi-classe (12 classes),
fenêtres partielles, \texttt{transitions\_only} + fin de fenêtre
& $\approx$ 0,59--0,64 & Vs baseline majorité $\approx$0,04 \\
\hline
\end{tabularx}
\caption{Anticipation — résultats}
\label{tab:anticipation}
\end{table}

La prédiction multi-classe de l'activité suivante depuis un contexte partiel
atteint un macro-F1 d'environ 0,60--0,64 selon $p$.

\section{Contributions supplémentaires}

\begin{table}[H]
\centering
\small
\begin{tabularx}{0.75\textwidth}{|Y|c|}
\hline
\textbf{Module} & \textbf{Métrique principale} \\
\hline
Monte Carlo Dropout ($T=20$) \parencite{gal2016} & Entropie épistémique pour UWR \\
Fall-risk binaire (transitions 9--12)            & Macro-F1 $\approx$ 0,88 \\
Détection impact (hybride historique)            & AUC-ROC = 0,762 \\
Adaptateur CORAL \parencite{sun2016}             & F1 cross-dataset = 0,629 \\
SSL SimCLR-style / foundation-style              & Checkpoints \texttt{ssl\_backbone.pt}, \texttt{foundation\_style.pt} \\
Calibration + bandit de seuil                    & Modules online (démo Streamlit) \\
\hline
\end{tabularx}
\caption{Synthèse des modules auxiliaires}
\label{tab:aux_modules}
\end{table}

\subsection{Fall-risk, SSL et calibration}

Sur le protocole anticipation binaire fall-risk
(\texttt{train\_fall\_risk.py}), le macro-F1 de validation atteint 0,88 pour
$p\in\{50\%,75\%\}$ (accuracy $\approx$0,925). Le SSL et le pipeline
foundation-style produisent des checkpoints utilisables comme initialisations
optionnelles ; ils ne remplacent pas le pré-entraînement supervisé de
production dans les tableaux CL principaux. La calibration de température et
le bandit de seuil sont validés fonctionnellement dans l'interface Streamlit
(adaptation en ligne du seuil d'alerte).

\section{Guide de reproductibilité}

Pour reproduire les résultats du chapitre~6 à partir du dépôt \texttt{har\_project/} :

\begin{enumerate}
  \item Installer les dépendances : \texttt{pip install -r requirements.txt}.
  \item Placer HAPT, WISDM et PAMAP2 dans \texttt{data/raw/} (cf. annexe~A).
  \item Lancer le prétraitement : \texttt{python scripts/preprocess.py}.
  \item Pré-entraînement : \texttt{python scripts/train.py --mode pretrain}.
  \item CL user-incrémental :
  \texttt{python scripts/train.py --mode continual --scenario user}.
  \item CL class-incrémental :
  \texttt{python scripts/train.py --mode continual --scenario class}.
  \item Anticipation :
  \texttt{python scripts/train\_anticipation.py}.
  \item Figures : \texttt{bash scripts/regenerate\_figures.sh}.
\end{enumerate}

Les hyperparamètres par défaut correspondent aux tableaux annexe~B. Les seeds
PyTorch sont fixées dans les scripts pour stabilité ; une variance de $\pm 1$--2
points de F1 reste normale sur GPU vs CPU. Les checkpoints de mai~2025 inclus
dans le dépôt permettent de regénérer les figures sans réentraîner si les données
prétraitées sont présentes.

\section{Discussion générale}

\subsection{Synthèse}

\begin{table}[H]
\centering
\small
\begin{tabularx}{0.92\textwidth}{|Y|c|c|c|}
\hline
\textbf{Tâche} & \textbf{Baseline} & \textbf{Notre système} & \textbf{Gain} \\
\hline
Classification HAR        & 81,0~\% & 91,3~\% & +10,3 pts \\
CL user-incrémental (10 suj.) & 62,4~\% & 79,2~\% & +16,8 pts \\
CL class-incrémental        & 58,1~\% & 74,6~\% & +16,5 pts \\
Anticipation LSTM multi-classe & — & $\approx$0,60--0,64 & fin fenêtre + poids \\
Cross-dataset CORAL         & 0,040   & 0,629   & $\times 15$ \\
Fall-risk (macro-F1 binaire) & — & $\approx$0,88 & transitions 9--12 \\
\hline
\end{tabularx}
\caption{Bilan quantitatif global}
\label{tab:bilan_global}
\end{table}

\subsection{Limites}

\begin{itemize}
  \item \textbf{Class-incrémental} : F1 final modeste (0,16--0,17) — déséquilibre
  extrême entre classes rares et dominantes.
  \item \textbf{Cross-domain} : oubli marqué à l'arrivée de PAMAP2 (tâches 37--44).
  \item \textbf{Anticipation} : macro-F1 multi-classe d'environ 0,60--0,64 selon $p$.
  \item \textbf{Chutes} : proxy HAPT (fall-risk transitions), pas de chutes réelles
  (MobiAct en perspective) ; repas non entraînable sans corpus cuisine.
  \item \textbf{SSL / foundation / RL} : modules légers démonstratifs, pas un FM
  public ni du deep RL.
\end{itemize}

\subsection{Perspectives d'amélioration}

Trois axes ressortent des expériences pour un déploiement réel.

\textbf{Tampon domain-aware.} L'oubli PAMAP2 (tâches 37--44) suggère d'étiqueter
les exemples du rejeu par domaine capteur et de garantir un quota minimum lors
de l'échantillonnage UWR. CORAL intermédiaire avant chaque nouveau domaine
pourrait stabiliser les embeddings.

\textbf{Anticipation.} Améliorer le contexte temporel et le déséquilibre
(oversampling des transitions, entraînement conjoint plus poussé) pour relever
le macro-F1 au-delà de $\approx$0,64.

\textbf{Validation chutes et ADLs.} Intégrer MobiAct pour les chutes réelles et
un corpus cuisine (Opportunity / Cooking) pour le scénario repas clôtureraient
les exemples d'application de la fiche.

\section{Conclusion}

Les expériences confirment la chaîne conception (ch.~4) $\rightarrow$ implémentation
(ch.~5) : backbone compétitif, UWR efficace, anticipation LSTM opérationnelle
mais encore limitée, CORAL synergique avec le pré-entraînement. Les axes
d'amélioration (tampon plus grand, MobiAct, anticipation) sont identifiés
pour la suite.
